\begin{flushright}
  \fechaclase{08 de septiembre de 2026}
\end{flushright}

Haciendo
\[
  \psi=PX
\]
\[
  \dot{V}
  =\psi^{T}AX+X^{T}A^{T}\psi
  -\psi^{T}BF\psi-\psi^{T}F^{T}B^{T}\psi
\]
como
\[
  X=P^{-1}PX=P^{-1}\psi
\]
y dado que $P$ es no singular decimos que
\[
  P^{-1}=\gamma
\]
entonces
\[
  X=\gamma\psi
\]
Así,
\[
  \dot{V}
  =\psi^{T}A\gamma\psi+\psi^{T}\gamma A^{T}\psi
  -(\psi^{T}BF\psi+\psi^{T}F^{T}B^{T}\psi)
\]

Los cambios de variable convierten el problema en una LMI, para que el
sistema sea globalmente asintóticamente estable
\[
  \dot{V}
  =\psi^{T}\left[A\gamma+\gamma A^{T}-(BF+F^{T}B^{T})\right]\psi<0
\]

Se necesita que
\[
  A\gamma+\gamma A^{T}-(BF+F^{T}B^{T})<0
\]
con el complemento de Schur
\[
  \begin{bmatrix}
    A\gamma+\gamma A^{T}-(BF+F^{T}B^{T}) & 0\\
    0 & -I
  \end{bmatrix}<0
\]
\[
  \gamma>0
\]
\[
  P=\gamma^{-1}
\]
\[
  K=FP
\]

Las ganancias se calculan al resolver la LMI anterior.

\paragraph{\tituloejemplo{Ejemplo}}
Diseñe un controlador por realimentación de estados para el sistema
\[
  \dot{X}_1=X_2
\]
\[
  \dot{X}_2=2X_1+3X_2+v
\]
\[
  \begin{bmatrix}
    \dot{X}_1\\
    \dot{X}_2
  \end{bmatrix}
  =
  \begin{bmatrix}
    0 & 1\\
    2 & 3
  \end{bmatrix}
  \begin{bmatrix}
    X_1\\
    X_2
  \end{bmatrix}
  +
  \begin{bmatrix}
    0\\
    1
  \end{bmatrix}v
\]
